\chapter{Conclusões}\label{chap:conclusoes}
As conclusões devem responder aos objetivos, sintetizar as contribuições sustentadas pelos resultados e reconhecer limitações. Este texto é um exemplo de organização e deve ser substituído pelas conclusões efetivas do trabalho.

\section{Síntese e contribuições}
Recupere a pergunta de pesquisa e indique, de maneira proporcional às evidências, o que foi demonstrado, desenvolvido ou observado. Evite acrescentar resultados inéditos nesta seção.

\section{Perspectivas futuras}
Apresente extensões plausíveis, questões abertas e possibilidades de aprofundamento. Nos modos de qualificação, as atividades futuras já estão organizadas também no cronograma.
